\section{实验结果与性能分析}

\subsection{FID 对比}

讲者展示了不同采样方法在 CIFAR-10 等数据集上的 FID（Frechet Inception Distance）对比：

\begin{itemize}
    \item \textbf{DDPM}（1000 步）：FID 约 3$\sim$4（基线）
    \item \textbf{DDIM}（100 步）：FID 有明显退化，约 5$\sim$8
    \item \textbf{DDIM}（10 步）：FID 大幅恶化，约 20+
    \item \textbf{DPM-Solver-2}（15$\sim$20 步）：FID 接近 DDPM 1000 步的水平
    \item \textbf{DPM-Solver-3}（10$\sim$15 步）：FID 达到甚至优于 DDPM 1000 步
\end{itemize}

\begin{importantbox}{关键结果}
使用 DPM-Solver-3，仅需 15 步即可达到 DDPM 1000 步的采样质量。这意味着约 \textbf{70 倍}的加速，且\textbf{无需重新训练模型}——完全在推理阶段通过更优的 ODE 求解器实现。这也是讲者反复强调的一个亮点：不同的数学视角（SDE $\to$ ODE $\to$ 半线性结构 $\to$ 指数积分器）可以直接转化为实际性能提升。
\end{importantbox}

\subsection{不同阶次的收敛行为}

\begin{knowledgebox}{阶次与步数的关系}
\begin{itemize}
    \item 步数充足（$>50$）时，各阶 DPM-Solver 差异不大，因为欧拉法本身误差已经很小
    \item 步数较少（$10\sim 20$）时，高阶方法优势显著，二阶比一阶好一个数量级
    \item 步数极少（$<10$）时，三阶比二阶仍有提升，但收益递减
    \item 超过三阶后，网络本身的近似误差（$\boldsymbol{\epsilon}_\theta \neq \boldsymbol{\epsilon}^*$）成为瓶颈
\end{itemize}
\end{knowledgebox}

\subsection{与其他加速方法的关系}

DPM-Solver 属于\textbf{Training-free}（无需额外训练）的采样加速方法。与之互补的还有：

\begin{table}[H]
\centering
\begin{tabular}{lcc}
\toprule
\textbf{方法} & \textbf{需要训练？} & \textbf{核心思路} \\
\midrule
DPM-Solver/DDIM & 否 & 更好的 ODE 离散化 \\
知识蒸馏（Progressive Distillation） & 是 & 学生模型模仿多步教师 \\
一致性模型（Consistency Models） & 是 & 直接学习从噪声到数据的映射 \\
Flow Matching + Euler & 否 & 更直的轨迹，欧拉法更有效 \\
\bottomrule
\end{tabular}
\caption{不同采样加速策略的比较}
\end{table}

\begin{knowledgebox}{DPM-Solver 与 Flow Matching 的互补}
Flow Matching 通过训练使 ODE 轨迹更直（接近线性），从而让欧拉法就能工作得很好。DPM-Solver 则是在给定（可能弯曲的）ODE 轨迹上寻找更好的数值求解方案。两者的思路正交且互补：如果 ODE 轨迹已经很直，DPM-Solver 的优势会减小；但对于标准扩散模型的弯曲轨迹，DPM-Solver 的价值就非常显著。
\end{knowledgebox}

\subsection{本章小结}

实验证明，DPM-Solver 在不重新训练模型的前提下，将采样步数从 1000 步压缩到 10$\sim$20 步，实现了数十倍的加速。其核心优势在于利用了扩散模型 ODE 的半线性结构，这是通用 ODE 求解器无法做到的。


